\chapter{Variáveis e Escopos}

\section{Declaração com \texttt{let}}

Toda variável é declarada com \lstinline{let}. A declaração vincula um nome a um valor:

\begin{lstlisting}
let x = 10
let nome = "João"
let ativo = true
\end{lstlisting}

\lstinline{let} sem atribuição inicial não existe em \hayashi{} — toda variável é
inicializada no momento da declaração.

\subsection*{Redeclaração}

Uma variável pode ser redeclarada com \lstinline{let} no mesmo escopo. O novo valor
substitui o anterior sem afetar outras variáveis que apontem para o valor antigo:

\begin{lstlisting}
let x = 5
let y = x      // y aponta para 5
let x = 99     // x agora é 99; y continua 5
display y      // 5
\end{lstlisting}

\section{Constantes com \texttt{const}}

\lstinline{const} declara um nome que não pode ser redeclarado nem atribuído no
mesmo escopo. Útil para valores de configuração ou parâmetros fixos:

\begin{lstlisting}
const ALPHA = 0.05
const NOME_PROJETO = "LAFI"
\end{lstlisting}

Tentar atribuir a uma constante produz erro de compilação:

\begin{lstlisting}
const N = 100
N = 200    // error: cannot assign to constant 'N'
\end{lstlisting}

\begin{nota}
  \lstinline{const} em \hayashi{} é avaliado em tempo de execução, não em tempo de
  compilação. A distinção de \lstinline{let} é puramente de mutabilidade de ligação,
  não de desempenho.
\end{nota}

\section{Atribuição}

Após declaração com \lstinline{let}, a variável pode ser reatribuída com \lstinline{=}:

\begin{lstlisting}
let contador = 0
contador = contador + 1
contador = contador + 1
display contador    // 2
\end{lstlisting}

Operadores compostos de atribuição:

\begin{lstlisting}
let n = 10
n += 5     // n = 15
n -= 3     // n = 12
n *= 2     // n = 24
n /= 4     // n = 6
\end{lstlisting}

\section{Escopos}

Escopos são delimitados por blocos \lstinline|\{ \}|. Variáveis declaradas dentro de
um bloco existem apenas até o final dele:

\begin{lstlisting}[caption={Variáveis locais a um bloco}]
let x = 1

{
  let y = 2
  display x    // 1  — x é visível aqui
  display y    // 2
}

display x    // 1
display y    // error: undefined variable 'y'
\end{lstlisting}

\subsection*{Escopos em funções}

Cada chamada de função cria um novo escopo. Parâmetros e variáveis locais são
destruídos quando a função retorna:

\begin{lstlisting}
fn soma(a, b) {
  let resultado = a + b
  return resultado
}

display soma(3, 4)    // 7
display resultado     // error: undefined variable 'resultado'
\end{lstlisting}

\subsection*{Captura de escopo exterior (closures)}

Funções anônimas (closures) capturam variáveis do escopo onde foram definidas:

\begin{lstlisting}
let base = 100

let adicionar = fn(x) {
  return x + base    // captura 'base' do escopo exterior
}

display adicionar(20)    // 120
\end{lstlisting}

\section{Variáveis especiais}

\subsection*{\texttt{\_n} e \texttt{\_N} em contextos de dados}

Dentro de \lstinline{generate} e \lstinline{mutate}, duas variáveis implícitas
estão disponíveis:

\begin{itemize}
  \item \lstinline{_n}: índice da linha atual (começando em 1)
  \item \lstinline{_N}: número total de linhas do DataFrame
\end{itemize}

\begin{lstlisting}[caption={Uso de \_n e \_N}]
input x
  10  20  30
end

generate df id = _n       // [1, 2, 3]
generate df peso = 1/_N   // [0.333, 0.333, 0.333]
\end{lstlisting}

\subsection*{\texttt{\_} (descarte)}

O identificador \lstinline{_} descarta um valor sem criar uma ligação:

\begin{lstlisting}
for _ in 1..5 {
  display "iteração"
}
\end{lstlisting}

\section{Blocos expressão}

Um bloco pode ser usado como expressão. Ele avalia uma sequência de statements e retorna o valor da última expressão. Variáveis declaradas dentro do bloco são locais a ele:

\begin{lstlisting}[caption={Bloco expressão retornando um valor}]
let df = {
  let raw = load("dados.csv")
  generate raw y = log(x)
  keep(raw, ["date", "y"])
  raw
}
\end{lstlisting}

Após o bloco, \lstinline{df} está disponível, mas \lstinline{raw} não.

\section{Controle de saída}

\lstinline{quietly on} suprime a saída automática de statements e estimadores. \lstinline{print(...)} e \lstinline{display ...} continuam aparecendo. \lstinline{quietly off} restaura a saída normal. A flag respeita escopo:

\begin{lstlisting}[caption={quietly on/off com escopo}]
quietly on

let df = {
  let a = load("a.csv")
  let b = load("b.csv")
  let m = merge(a, b, key=id, type=inner)
  generate m z = x - y
  m
}

quietly off

ols(z ~ x, df)
print("pronto")
\end{lstlisting}

Dentro do bloco, \lstinline{quietly off} habilitaria temporariamente a saída; ao sair do bloco, ela volta para \lstinline{quietly on}.
